\chapter{Curriculum Vitæ}
\label{AppendixB}
\vspace{-1cm}
Helena Rasche was born on March 22, 1992, in Bærum, Oslo, Norway. She lived all around the world, completing primary school in Oklahoma in the US, before middle school in Singapore, and high school in Texas. It was then, in high school, that she took her one and only formal course in computer science, introductory Java.

She began her university studies in Chemical Engineering at Texas A\&M University. She took a few biology and genomics classes at that time, considering that she might want to later become a medical doctor, but found it so much more engaging than Chemical Engineering that she switched her major and completed her BSc in Biochemistry in 2013. All throughout this time she was increasingly passionate about this new thing called Linux and system administration, even teaching herself Perl. After graduation she began working with the Bacteriophage research group at Texas A\&M where she got to apply all of her new software development and system administration knowledge to help them update their infrastructure and design new more reproducible analyses. In her youthful enthusiasm she completely rewrote a large portion of the course materials for their undergraduate annotation course to use the CPT's newly developed phage annotation pipelines. She worked with the group to release all of their software as open source, benefiting the entire world of phage annotators.

She moved next to Freiburg, Germany where she oversaw the second largest deployment of Galaxy in the world, designing many of the scaling best practices that are still in use today, scaling their services from 500 users to 10.000 when she left. On the side she helped administer the BWCloud Freiburg OpenStack deployment.

Finally she moved to the Netherlands and found her place at Erasmus MC where she worked on various projects in the Department of Pathology, with a special focus on reproducibility and teaching. She became a co-lead for the Galaxy Training Network, combining her passion for teaching and science and reproducible, open research. Her incredible partner Saskia and she co-wrote an Erasmus+ grant which funded a large portion of her time at ErasmusMC. Thanks to this grant, after years spent scaling infrastructure, she now had experience scaling live trainings events from 20 to 2.000 attendees.

She has gone on to a postdoc position in the Clinical Genetics department at nearby LUMC.
